%!TEX encoding = UTF-8 Unicode
%!TEX program = xelatex

\documentclass[doctor]{hnuthesis}
\usepackage[small,bf,sf,SF,BF]{subfigure}
\usepackage{amsmath, unicode-math}
\DeclareMathAlphabet{\mathcal}{OMS}{cmsy}{m}{n}
\DeclareMathAlphabet{\mathbb}{U}{msb}{m}{n}%
%\setlength\cftbeforechapskip{0.5em}  
%\newCJKfontfamily{\simsun}[AutoFakeBold = {3.17}]{SimSun} %%中文加粗
\setlength{\algoheightrule}{1.5pt} % thickness of the rules above and below
\setlength{\algotitleheightrule}{0.75pt} % thicknes of the rule below the title
\usepackage{algpseudocode}
 
 
\usepackage{float}
\usepackage{makecell}





\usepackage{tikz}
\usepackage{etoolbox}
\newcommand{\circled}[2][]{\tikz[baseline=(char.base)]
    {\node[shape = circle, draw, inner sep = 1pt]
    (char) {\phantom{\ifblank{#1}{#2}{#1}}};%
    \node at (char.center) {\makebox[0pt][c]{#2}};}}
\robustify{\circled} 

\usepackage{listings}

%\usepackage{setspace}
%\linespread{1.5}



% 学校代码
\hnucode{10532}
% 学校名称
\hnuname{湖南大学}
\enhnuname{Hunan University}
% 中图分类号
\clc{TP391}         
% 密级       
\secrettext{普通}                  

% 标题 
\title{智能交通系统的决策控制方法研究}
\entitle{Research on decision control methods for intelligent transportation systems}
% 作者
\author{XXXXXX}
\enauthor{XXXXXX}
% 学号
\authorid{***********}      
% 学院
\college{机械与运载工程学院}
% 专业
\major{机械工程} 
\enmajor{Mechanical Engineering}
%学士学位获得学校，年份
\enbachelor{B.E.~(****** University)2016}
\enmaster{M.S.~(****** University)2020}
\endoctor{Doctor of Philosophy}
% 研究方向
\workon{智能网联汽车}
% 导师
\supervisor{***\ 教授}
\ensupervisor{Professor XXX}
% 论文提交、答辩日期
\submitdate{202x年x月xx日}
\defensedate{202X年x月xx日}
\endate{June, 2022}
% 答辩委员会主席
\chair{待定}

\usepackage[export]{adjustbox} 
\begin{document}
% 封面、原创性声明
\maketitle

% 摘要
%中文摘要
\begin{abstract} %TODO摘要没有居中
	中文摘要。文章
	\keywords{关键字1；关键字2；关键字n}
\end{abstract}

%英文摘要
\begin{enabstract}
	英文摘要。
	\enkeywords{Keyword1; Keyword2; Keywordn}
\end{enabstract}
% 目录
\tableofcontents
% 插图附表索引
\begingroup
	\renewcommand*{\addvspace}[1]{}
	\newcommand{\loflabel}{图} 
	\renewcommand{\numberline}[1]{\loflabel~#1\hspace*{1em}}	
	\listoffigures
	
	\newcommand{\lotlabel}{表}
	\renewcommand{\numberline}[1]{\lotlabel~#1\hspace*{1em}}
	\listoftables
\endgroup


% 正文章节
\mainmatter
\chapter{引\quad 言}

绪论。

\begin{figure}
    \centering    \includegraphics{figures/hnu-logo.png}
    \caption{Caption标注}
    \label{fig:1}
\end{figure}


\chapter{论文模板使用说明}
\section{Latex使用}
该模板已经在mac的texifier，overleaf，vscode上验证使用，请用XeLaTex进行编译。

%\section{数学公式、图表、参考文献等功能}


\section{参考文献}

在Doctoral-thesis.tex修改参考文献链接：\lstinline!\bibliography{MyLibrary.bib}!。

目前可以使用的引用命令有\lstinline!\cite{}!、\lstinline!\citep{}!、\lstinline!\citet{}!和\lstinline!\citen{}!。

效果如下：
\begin{enumerate}
	\item 引用参考文献时，数字显示在右上角：引用 ResNet~\cite{2020_aradi_Survey}提出了ResNet网络结构\citep{2020_aradi_Survey}（\lstinline!\cite{}!、\lstinline!\citep{}!）；
	\item 引用参考文献时，以人名为主语：\citet{2017_qinxiaohui_FeiYunZhiCheLiangDuiLieDeFenBuShiKongZhi}针对队列稳定性提出了观测性方法。或者\citet{2021_chen_Graph}对与图做了个总结（\lstinline!\citet{}!）；
	\item 引用参考文献时，引用数字显示在正文中：\citen{2017_qinxiaohui_FeiYunZhiCheLiangDuiLieDeFenBuShiKongZhi}这样（\lstinline!\citen{}!）。
\end{enumerate}
%引用 ResNet~\cite{2020_aradi_Survey}提出了ResNet网络结构\citep{2020_aradi_Survey}（上述第一个效果和第二个效果相同）。

%以人名为主语：\citet{2017_qinxiaohui_FeiYunZhiCheLiangDuiLieDeFenBuShiKongZhi}针对队列稳定性提出了观测性方法。或者\citet{2021_chen_Graph}对与图做了个总结。

%参考文献的引用图标不上标：\citen{2017_qinxiaohui_FeiYunZhiCheLiangDuiLieDeFenBuShiKongZhi}这样。

\section{图表}
插入图的常用命令有：

插入单个图时，可使用以下命令，主要调整的是图片大小。
\begin{lstlisting}
\begin{figure}[htb]
  \centering
  \includegraphics[width=.7\textwidth]{figures/hnu-logo.png}
  \caption{湖南大学校徽}
  \label{fig:ch2-1}
\end{figure}
\end{lstlisting}

引用时可以直接使用\lstinline!\figref{fig:ch2-1}!，文本中显示为“\figref{fig:ch2-1}”。值得一提的是，如果插入的为pdf文件，其中每页pdf对应着一张图片，那么，\lstinline!\includegraphics[]!中可添加\lstinline!page=num!（num为对应的页码）。此外需要提醒的是虽然\lstinline!width=!可以调整图片的大小，即可手动放大缩小来适应排版，但是\textcolor{red}{如果图中有文字的不太建议这样调整，最好是提前测试图片中的文字大小，不可小于8号字体，最大也不要超过10号字体，然后插入的时候尽量不使用\lstinline!width=!来调整，以保持原图的大小。}
\begin{figure}
    \centering
    \includegraphics[]{figures/hnu-logo.png}
    \caption{湖南大学校徽}
    \label{fig:ch2-1}
\end{figure}

多图排列，可使用以下代码：
\begin{lstlisting}
\begin{figure}[htbp]
  \centering
  \subfigure[分流]{
  \includegraphics[width=0.3\linewidth]{figures/hnu-logo.png}
  \label{fig:ch2-2a}
  }
  \quad
  \subfigure[交叉]{
  \includegraphics[width=0.3\linewidth]{figures/hnu-logo.png}
  \label{fig:ch2-2b}
  }
  \quad
  \subfigure[合流]{
  \includegraphics[width=0.3\linewidth]{figures/hnu-logo.png}
  \label{fig:ch2-2c}
  }
  \quad
  \subfigure[无冲突]{
  \includegraphics[width=0.3\linewidth]{figures/hnu-logo.png}
  \label{fig:ch2-2d}
  }
  \caption{冲突模式}
  \label{fig:ch2-2}
\end{figure}
\end{lstlisting}
\begin{figure}[htbp]
  \centering
  \subfigure[分流]{
  \includegraphics[width=0.3\linewidth]{figures/hnu-logo.png}
  \label{fig:ch2-2a}
  }
  \quad
  \subfigure[交叉]{
  \includegraphics[width=0.3\linewidth]{figures/hnu-logo.png}
  \label{fig:ch2-2b}
  }
  \quad
  \subfigure[合流]{
  \includegraphics[width=0.3\linewidth]{figures/hnu-logo.png}
  \label{fig:ch2-2c}
  }
  \quad
  \subfigure[无冲突]{
  \includegraphics[width=0.3\linewidth]{figures/hnu-logo.png}
  \label{fig:ch2-2d}
  }
  \caption{冲突模式}
  \label{fig:ch2-2}
\end{figure}


~\\
~\\
三线表的命令：
\begin{lstlisting}
\begin{table}[htb]
  \centering
  \caption{训练超参数}
  \label{tab:ch2-1}
  \begin{tabular}{ccc} %ccc表示列数与对齐方式
	\toprule[1.5pt] % 标准尺寸不要修改
	物理意义 & 英文/符号 & 数值 \\ % 每一列的元素用&隔开
	\midrule[0.75pt]
	 隐含层数 & Hidden layer number & $\rm 2$ \\
	\bottomrule[1.5pt]
  \end{tabular}
\end{table}
\end{lstlisting}

引用格式为\lstinline!\tabref{tab:ch2-1}!，在文本中为“\tabref{tab:ch2-1}”。
\begin{table}[htb]
%  \small
  \centering
  \caption{单车强化学习MDP设置}
  \label{tab:ch2-1}
  \begin{tabular}{cc}
	\toprule[1.5pt]
	名称 & 参数  \\ 
	\midrule[0.75pt]
	状态 & \makecell{鸟瞰图 \\ 显示航向角和速度的图像快照 \\ 车辆特征向量和信号特定信息} \\
	动作 &  \makecell{加速度 \\ 等待、慢慢前进、走 \\ 向上、向下、向左和向右移动} \\
	奖励 & \makecell{成功到达奖励和碰撞惩罚 \\ 速度、跟车间距和预计到达交叉路口时间的函数} \\
	\bottomrule[1.5pt]
  \end{tabular}
\end{table}

\section{定理类环境}
该模板定义了定理，引理，评注等，具体有：
\begin{lstlisting}
\newcommand\hnu@assertionname{断言}
\newcommand\hnu@axiomname{公理}
\newcommand\hnu@corollaryname{推论}
\newcommand\hnu@definitionname{定义}
\newcommand\hnu@propertyname{性质}
\newcommand\hnu@examplename{例}
\newcommand\hnu@lemmaname{引理}
\newcommand\hnu@proofname{证明}
\newcommand\hnu@propositionname{命题}
\newcommand\hnu@remarkname{评注}
\newcommand\hnu@theoremname{定理}
\newcommand\hnu@assumptionname{假设}
\end{lstlisting}

\begin{lstlisting}
\begin{proof}
	这是一个证明。
\end{proof}

\begin{theorem}\label{the:ch2-1}
	这是一个定理。
\end{theorem}
\end{lstlisting}

\begin{proof}
	这是一个证明。
\end{proof}

\begin{theorem}\label{the:ch2-1}
	这是一个定理。
\end{theorem}

\section{公式书写与引用}
公式书写的开始形式如下
引用格式为\lstinline!\eqref{}!


\section{伪代码书写}
\begin{algorithm} [htb]                                            
	\caption{FCFS算法}
  	\label{algo:revi-1}  
  	\SetKwInOut{KwIn}{初始化}
    \SetKwInOut{KwOut}{输出}
	\KwIn{交叉路口信息，即网格占用点}
	得到交叉路口车辆进入序列 \\
	\For{\rm $1 < j < N$}{
		\For{\rm $k<j$}{
			识别车辆$j$与$k$之间的公共占用点$c_{j,k}\in \{1,2,\dots,Q_j\}$ \\
			获得车辆$k$的进入时间$t_k^{\rm in}$ \\
			$t_{j, k}^{\rm in}=t_k^Q+\frac{S}{V_{\text {int}}}+\delta$ \quad $\rhd$ 假设$Q$是占用点
			}
		$t_j^{\mathrm{in}}=\max \left\{t_{j, 1}^{\mathrm{in}}, t_{j, 2}^{\mathrm{in}}, \ldots, t_{j, k}^{\mathrm{in}}\right\}$ \quad $\rhd$ 获得最保守的进入时间 \\
		随后更新并存储$t_j^{\mathrm{in}}, t_j^1, t_j^2, \ldots t_j^{Q_j}, t_j^{\text {out }}$。
		}
\end{algorithm}

\begin{summary}
	总结。引用 ResNet~\citep{2020_aradi_Survey}提出\citen{2020_aradi_Survey}。\citep{2021_chen_Graph, 2021_qiao_Reinforcement, 2021_di_survey, 2017_qinxiaohui_FeiYunZhiCheLiangDuiLieDeFenBuShiKongZhi}中文显示\citen{2017_qinxiaohui_FeiYunZhiCheLiangDuiLieDeFenBuShiKongZhi}\citen{2020_aradi_Survey}效果后\citet{2017_qinxiaohui_FeiYunZhiCheLiangDuiLieDeFenBuShiKongZhi}面的问下好的看看\citet{2021_chen_Graph}提出了\citen{2017_qinxiaohui_FeiYunZhiCheLiangDuiLieDeFenBuShiKongZhi}可行吗？
\end{summary}

\bibliography{MyLibrary.bib}

\appendix
\chapter{攻读学位期间所发表的学术论文}

\begin{enumerate}
    \item 第三类永动机
\end{enumerate}


\chapter{攻读学位期间所参加的科研项目}

\begin{enumerate}
    \item 第三类永动机
\end{enumerate}


\backmatter
\begin{acknowledgements}

	致谢。

\end{acknowledgements}


\end{document}
